\documentclass[11pt]{article}

% ---- Packages ----
\usepackage[margin=1in]{geometry}
\usepackage{graphicx}      % figures
\usepackage{booktabs}      % table lines from kable (booktabs = TRUE)
\usepackage{float}         % [H] keeps figures where they are written
\usepackage[font=small,labelfont=bf]{caption}
\usepackage{hyperref}


\title{Maternal Vitamin D, Obesity and Immune Gene Expression in Preeclampsia}
\author{Irene Chen \\ QBS 103 Final Project}
\date{Fall 2026}

\begin{document}
\maketitle

% =====================================================================
\section{Methods}

\paragraph{Data source.}
Data came from a nested case-control study within the Vitamin D Antenatal Asthma Reduction
Trial (VDAART) \cite{mirzakhani2016}, available in the Gene Expression Omnibus as GSE85307.
Peripheral whole-blood gene expression was measured at 10--18 weeks of gestation with the
Affymetrix Human Gene 1.0 ST Array; values are on the log$_2$ scale (range 3.1--13.2). The
initial data set contained 157 samples and 500 probes. Covariates were
matched to the expression data by GEO sample ID, and all 157 samples matched. Probe IDs were
linked to gene symbols with the provided linker file.

\paragraph{Final sample size.}
No samples had missing values for the variables analyzed, so no data points were removed. The
final sample size was 157 women: 110 with normal pregnancies and 47 who developed preeclampsia.

\paragraph{Variables.}
The genes of interest were \textit{LYN}, \textit{JAK2} and \textit{PTX3}, chosen because
preeclampsia involves maternal immune and inflammatory changes. The continuous covariate of
interest was maternal 25-hydroxyvitamin D (ng/mL). The categorical covariates were pregnancy
outcome (normal vs.\ preeclampsia) and maternal obesity status (not obese, BMI $<$ 30 kg/m$^2$;
obese, BMI $\geq$ 30 kg/m$^2$), derived from continuous BMI. Maternal age (years) and
gestational age at delivery (weeks) were included as two additional continuous covariates in
the summary table.

\paragraph{Analyses.}
Continuous covariates were summarized as mean (standard deviation) and categorical covariates
as n (\%), stratified by pregnancy outcome. The distribution of \textit{LYN} expression was
shown with a histogram, its association with vitamin D was measured with a Pearson correlation
and a linear fit, and its expression by pregnancy outcome and obesity status was shown with
boxplots. Ten genes (\textit{LYN}, \textit{JAK2}, \textit{PTX3} and the seven genes with the
highest standard deviation of expression across all women among the 500 probes) were shown in a heatmap of
row-scaled (z-score) expression, with hierarchical clustering of genes and samples (Euclidean
distance, complete linkage). Finally, a dumbbell plot compared the mean z-score of the same ten
genes between pregnancy outcomes. All analyses were descriptive.

\paragraph{Software.}
All analyses were performed in R version 4.5.3 \cite{R2026} with the packages ggplot2
(v4.0.3) \cite{ggplot2}, ggpubr (v1.0.0) \cite{ggpubr}, pheatmap (v1.0.13) \cite{pheatmap}
and knitr (v1.52) \cite{knitr}.

\paragraph{Code availability.}
All code is available at \url{https://github.com/cheonijin927/QBS103-Final-Project}.

% =====================================================================
\section{Results}

\subsection{Participant characteristics}

Women who developed preeclampsia had lower mean vitamin D levels than women with normal
pregnancies (19.7 vs.\ 24.4 ng/mL), delivered earlier on average (37.8 vs.\ 39.3 weeks) and
were more often obese (55.3\% vs.\ 28.2\%), while maternal age was similar between groups
(25.8 vs.\ 26.6 years) (Table~\ref{tab:summary}).

% Table made in R and saved as table1.tex (upload it to Overleaf next to this file)
\input{table1.tex}

\subsection{\textit{LYN} expression}

\textit{LYN} expression was approximately symmetric around a mean of 12.18 log$_2$ units
(Figure~\ref{fig:lyn}A). \textit{LYN} expression was positively associated with maternal
vitamin D (Pearson $r$ = 0.39, $p$ $<$ 0.001; Figure~\ref{fig:lyn}B). Within each pregnancy
outcome, median \textit{LYN} expression was slightly lower in obese than in non-obese women
(normal: 12.05 vs.\ 12.26; preeclampsia: 12.23 vs.\ 12.30), but the differences were small
and the groups overlapped substantially (Figure~\ref{fig:lyn}C).

\begin{figure}[H]
  \centering
  \includegraphics[width=\textwidth]{figures/Figure1_LYN.pdf}
  \caption{\textit{LYN} expression in whole-blood samples from 157 pregnant women.
  (A) Distribution of \textit{LYN} expression (log$_2$ intensity); the dashed line marks the
  mean. (B) \textit{LYN} expression vs.\ maternal vitamin D (25-hydroxyvitamin D, ng/mL); each
  point is one woman, the line is the linear fit with its 95\% confidence band, and the Pearson
  correlation is shown in the upper left. (C) \textit{LYN} expression by pregnancy outcome and
  obesity status (not obese, BMI $<$ 30 kg/m$^2$; obese, BMI $\geq$ 30 kg/m$^2$). Boxes show
  the median and interquartile range, whiskers extend to 1.5 times the interquartile range, and
  points are individual women.}
  \label{fig:lyn}
\end{figure}

\subsection{Heatmap of ten genes}

The ten genes shown were \textit{LYN}, \textit{JAK2}, \textit{PTX3}, \textit{UTS2},
\textit{RUNDC3A}, \textit{TMEM176B}, \textit{CEACAM6}, \textit{CLIC2}, \textit{SLC1A5} and
\textit{ZNF594}. Clustering grouped genes by their shared expression pattern across women, while women
with preeclampsia and obese women were spread across the sample clusters rather than forming a
separate group (Figure~\ref{fig:heatmap}).

\begin{figure}[H]
  \centering
  \includegraphics[width=\textwidth]{figures/Figure2_heatmap.pdf}
  \caption{Heatmap of ten genes (\textit{LYN}, \textit{JAK2}, \textit{PTX3} and the seven most
  variable genes in the data set) in 157 pregnant women. Each column is one woman and each row
  is one gene. Colors show row-scaled expression (z-score of log$_2$ intensity; red is higher
  and blue is lower than that gene's mean). Genes and women were clustered by hierarchical
  clustering (Euclidean distance, complete linkage). Tracking bars show pregnancy outcome and
  obesity status (not obese, BMI $<$ 30 kg/m$^2$; obese, BMI $\geq$ 30 kg/m$^2$).}
  \label{fig:heatmap}
\end{figure}

\subsection{Mean expression by pregnancy outcome}

Compared with normal pregnancies, mean expression in preeclampsia was most increased for
\textit{UTS2} (difference in mean z-score, 0.21) and most decreased for \textit{RUNDC3A}
($-$0.38). All differences were smaller than 0.4 standard deviations, indicating only modest
differences between pregnancy outcomes for these genes (Figure~\ref{fig:dumbbell}).

\begin{figure}[H]
  \centering
  \includegraphics[width=0.8\textwidth]{figures/Figure3_dumbbell.pdf}
  \caption{Dumbbell plot of mean expression of the ten heatmap genes by pregnancy outcome.
  Expression of each gene was converted to a z-score across all 157 women; points show the mean
  z-score in normal pregnancies (green) and preeclampsia (pink), and the grey line shows the
  difference between groups. Genes are ordered from the largest increase (top) to the largest
  decrease (bottom) in preeclampsia.}
  \label{fig:dumbbell}
\end{figure}

% =====================================================================
\bibliographystyle{unsrt}
\bibliography{references}

\end{document}
